\documentclass{article}
\usepackage[utf8]{inputenc}
\usepackage{geometry}
\geometry{margin=1in}
\usepackage{booktabs}
\usepackage{tabularx}
\usepackage{xcolor}
\usepackage{colortbl}

% This command tells LaTeX to make the flexible 'X' column vertically centered (like an 'm' column) instead of top-aligned
\renewcommand{\tabularxcolumn}[1]{m{#1}}

\begin{document}

\begin{table}[htbp]
    \centering
    \renewcommand{\arraystretch}{1.5} 
    \caption{Security and usability differences between FIDO2 Passkey form factors.}
    \label{tab:passkey_form_factors}
    \begin{tabularx}{\textwidth}{
        >{\bfseries\raggedright\arraybackslash}m{2.8cm} % Vertically centered, bolded row headers
        >{\raggedright\arraybackslash}X               % Synced column
        >{\raggedright\arraybackslash}X               % Device-bound column
    }
        \toprule
        % Note: \textbf{} is omitted for "Authenticator" because the entire column is set to bold above
        Authenticator & \textbf{Synced} & \textbf{Device-bound} \\
        \midrule
        
        Platform & Created by the built-in platform authenticator. The credential can be discovered and used on multiple devices via cloud sync. & Created by a single device's platform authenticator. Only usable on that specific hardware device. \\
        
        \arrayrulecolor{gray!40}\midrule\arrayrulecolor{black} 
        
        Roaming & Created on a separately attached device, but the credential can be discovered and used on multiple devices. & Created on a separately attached device. The credential cannot leave the authenticator, but the physical authenticator can be moved to login on multiple devices. \\
        
        \bottomrule
    \end{tabularx}
\end{table}

\end{document}